
\documentclass{article}
\usepackage{colortbl}
\usepackage{makecell}
\usepackage{multirow}
\usepackage{supertabular}

\begin{document}

\newcounter{utterance}

\centering \large Interaction Transcript for game `wordle\_withclue', experiment `high\_frequency\_words\_clue\_no\_critic', episode 6 with Qwen3.5{-}4B.
\vspace{24pt}

{ \footnotesize  \setcounter{utterance}{1}
\setlength{\tabcolsep}{0pt}
\begin{supertabular}{c@{$\;$}|p{.15\linewidth}@{}p{.15\linewidth}p{.15\linewidth}p{.15\linewidth}p{.15\linewidth}p{.15\linewidth}}
    \# & \multicolumn{2}{c}{Player} && \multicolumn{2}{c}{Game Master} \\
    \hline

    \theutterance \stepcounter{utterance}  
    & & & \multicolumn{4}{p{0.6\linewidth}}{
        \cellcolor[rgb]{0.9,0.9,0.9}{
            \makecell[{{p{\linewidth}}}]{
                \texttt{\tiny{[P1$\langle$GM]}}
                \texttt{You are a language wizard who likes to guess words by using the given rules.} \\
\\ 
\texttt{Welcome to Wordle! You have six attempts to guess the target word, a valid English word of five lowercase letters (a{-}z). Please use the tags "explanation:" and "guess:" to provide a concise explanation for each guess.} \\
\\ 
\texttt{To help you make an informed guess, you will receive a clue for the word, such as} \\
\texttt{clue: snowy white.} \\
\\ 
\texttt{Here is an example guess based on the clue:} \\
\texttt{explanation: In the fairy tail Snow White, the girl is killed because she eats a poisoned apple. And the word apple has 5 letters.} \\
\texttt{guess: apple} \\
\\ 
\texttt{After each guess, your answer will be validated, and you will receive feedback indicating which letters are correct (green), which letters are correct but in the wrong position (yellow), and which letters are incorrect (red). This feedback can be useful in determining which letters to include or exclude in your next guess.} \\
\\ 
\texttt{For example, the feedback for "apple" might be:} \\
\texttt{guess\_feedback: a<yellow> p<yellow> p<green> l<yellow> e<red>} \\
\\ 
\texttt{The explanation should contain details about how both the guess\_feedback and clue work together in tandem to formulate a new guess.} \\
\\ 
\texttt{Do not generate any other text.} \\
\\ 
\texttt{Let's begin with your first guess.} \\
\\ 
\texttt{clue: tv's"\_\_\_dad"} \\
            }
        }
    }
     \\ \\

    \theutterance \stepcounter{utterance}  
    & \multicolumn{4}{p{0.6\linewidth}}{
        \cellcolor[rgb]{0.561,0.737,0.561}{
            \makecell[{{p{\linewidth}}}]{
                \texttt{\tiny{[P1$\rangle$GM]}}
                \texttt{explanation: The clue "tv's'\_\_\_dad'" suggests a word that fits into the phrase "TV's \_\_\_ dad," likely referring to a famous character or a common idiom. The word "grand" fits perfectly as in "TV's grand dad" (referring to a grandfather), but more commonly, the phrase might be referencing "TV's grandpa" or a specific show reference. However, considering the structure and common wordle words, "grand" is a strong candidate. Let's try "grand" first.} \\
\texttt{guess: grand} \\
            }
        }
    }
    & & \\ \\

    \theutterance \stepcounter{utterance}  
    & & & \multicolumn{2}{p{0.3\linewidth}}{
        \cellcolor[rgb]{0.9,0.9,0.9}{
            \makecell[{{p{\linewidth}}}]{
                \texttt{\tiny{[GM$|$GM]}}
                \texttt{attempts = 1} \\
\texttt{target\_word = major} \\
\texttt{target\_word\_clue = tv's"\_\_\_dad"} \\
\texttt{guess = grand} \\
\texttt{guess\_feedback = g<red> r<yellow> a<yellow> n<red> d<red>} \\
            }
        }
    }
    & & \\ \\

    \theutterance \stepcounter{utterance}  
    & & & \multicolumn{4}{p{0.6\linewidth}}{
        \cellcolor[rgb]{0.9,0.9,0.9}{
            \makecell[{{p{\linewidth}}}]{
                \texttt{\tiny{[P1$\langle$GM]}}
                \texttt{guess\_feedback: g<red> r<yellow> a<yellow> n<red> d<red>} \\
\\ 
\texttt{Provide your response only in this format.} \\
\texttt{explanation: details} \\
\texttt{guess: word} \\
            }
        }
    }
     \\ \\

    \theutterance \stepcounter{utterance}  
    & \multicolumn{4}{p{0.6\linewidth}}{
        \cellcolor[rgb]{0.561,0.737,0.561}{
            \makecell[{{p{\linewidth}}}]{
                \texttt{\tiny{[P1$\rangle$GM]}}
                \texttt{explanation: The feedback g<red> r<yellow> a<yellow> n<red> d<red> tells us that the word contains one 'r' and one 'a', but neither is in the second or fourth positions respectively. Since 'g' and 'n' and 'd' are incorrect, we need to find a five{-}letter word with an 'r' and an 'a' that fits the clue "TV's \_\_\_ dad". The phrase likely refers to "TV's grandpa" but "grandpa" has 6 letters. Re{-}evaluating the clue, it might be a specific reference like "TV's dad" where the missing word is "grand" but that's 4 letters. Wait, the clue is "tv's'\_\_\_dad". Perhaps it's "TV's *real* dad"? No, 'e' and 'l' aren't confirmed. Let's reconsider the phrase. Could it be "TV's *new* dad"? No. How about "TV's *old* dad"? No. Let's look at the letters again: r and a are present. Maybe the word is "rare"? No, 4 letters. "Radar"? No. "Arora"? No. Let's think of famous TV dads. "TV's *real* dad" {-}> real (r, e, a, l). We have r and a. If the word is "real", positions would be r{-}e{-}a{-}l. Feedback was r<yellow> a<yellow>. So r is not 2nd, a is not 4th. In "real", r is 1st, a is 3rd. This doesn't match the yellow positions if the guess was "grand" (g{-}r{-}a{-}n{-}d). In "grand", r was 2nd (yellow), a was 3rd (yellow). So r is not 2nd, a is not 3rd. The word must have r and a. Maybe "star"? s{-}t{-}a{-}r. r is 4th (red in guess), a is 3rd (yellow). If the word is "star", r should be green or yellow depending on position. Let's try "star".} \\
\texttt{guess: star} \\
            }
        }
    }
    & & \\ \\

    \theutterance \stepcounter{utterance}  
    & & & \multicolumn{2}{p{0.3\linewidth}}{
        \cellcolor[rgb]{0.9,0.9,0.9}{
            \makecell[{{p{\linewidth}}}]{
                \texttt{\tiny{[GM$|$GM]}}
                \texttt{Error: The length of the guessed word is not 5.} \\
            }
        }
    }
    & & \\ \\

    \theutterance \stepcounter{utterance}  
    & & & \multicolumn{2}{p{0.3\linewidth}}{
        \cellcolor[rgb]{0.9,0.9,0.9}{
            \makecell[{{p{\linewidth}}}]{
                \texttt{\tiny{[GM$|$GM]}}
                \texttt{game\_result = ABORT} \\
            }
        }
    }
    & & \\ \\

\end{supertabular}
}

\end{document}
